\chapter{Implementación}

En este capitulo se documenta la implementación del sistema multiagente,  dando continuidad a la Propuesta de
Solución presentada en el capítulo anterior. Mientras dicho capítulo describe el diseño conceptual del sistema, en esta sección se detalla la arquitectura construida, sus componentes y las decisiones técnicas adoptadas
durante el desarrollo.

El proceso de implementación exigió tomar decisiones técnicas, adaptar componentes frente a restricciones que solo se manifestaron duracte la contrucción del sistema. Entre las decisiones más relevantes adoptadas durante esta etapa se cuentan el remplazo del modelo de lenguaje especificado en la propuesta por uno efectivamente disponible para el proyecto, la incorporación de un motor de preprocesamiento orientado a formatos heteroéneos, la extensión de la memoria de largo plazo del sistema y el desarrollo de una capa de interfaz y persistencia que la propuesta original no contemplaba.

Cada una de estas decisiones se documenta a lo largo del capítulo junto con su motivación técnica, de modo de distinguir con claridad entre los elementos de la arquitectura que se mantuvieron sin alteraciones respecto del diseño original y aquellos que constituyen una desviación deliberada, justificada por las condiciones reales de desarrollo.


\section{Arquitectura Implementada}

La arquitectura implementada preserva, en su núcleo agéntico, la topología
descrita en la propuesta: los siete nodos que componen el grafo, cuatro
agentes basados en modelos de lenguaje y tres herramientas determinísticas, junto con las tres aristas condicionales que gobiernan su
ejecución, se construyeron sin alteraciones respecto del diseño original.
El Agente Triage, la herramienta de recuperación aumentada por generación,
el Agente Selector de Herramientas, la herramienta de consulta de
vulnerabilidades, el Agente Ingeniero de Mitigación, la herramienta de
validación de comandos y el Agente Auditor Final operan tal como fueron
concebidos en la propuesta, conservando las mismas responsabilidades,
relaciones de dependencia y el patrón Plan-and-Execute que sustenta su
diseño. Esta correspondencia entre lo diseñado y lo construido constituye
una validación temprana de la arquitectura propuesta en el Seminario de
Título: el ciclo de razonamiento agéntico definido en el capítulo anterior
demostró ser implementable sin necesidad de reformular su estructura
fundamental.

Lo que sí difiere de la propuesta es todo aquello que rodea a ese núcleo.
La implementación exigió resolver un conjunto de capas
adicionales necesarias para que el sistema operara como una propuesta de producto de
software completo: un motor de preprocesamiento con dos versiones activas,
capaces de interpretar formatos de log heterogéneos antes de que la
información llegue al grafo; una capa de persistencia y presentación, la cual se refleja en una interfaz web que contiene un módulo de historial de análisis y generación de reportes en formato PDF— que no formaba parte del alcance original del Capítulo 8; y ajustes de
configuración motivados por restricciones reales del proveedor del modelo
de lenguaje, que no podían anticiparse en la etapa de diseño. La Figura 9.1
presenta esta arquitectura real de forma integral, distinguiendo
visualmente el núcleo agéntico —heredado sin cambios de la propuesta— de
los componentes incorporados durante la implementación.

\begin{figure}[h!]
    \centering
    \includegraphics[width=0.95\textwidth]{fig_9_1_arquitectura_real.png}
    \caption{Arquitectura real de NODEX. El núcleo agéntico, al centro,
    corresponde a la Figura 8.3 de la propuesta; los componentes restantes
    (preprocesamiento con dos motores, capa de API, interfaz y ciclo de
    retroalimentación) se incorporaron durante la implementación. Fuente:
    elaboración propia.}
    \label{fig:9.1}
\end{figure}

En términos operativos, el sistema expone su funcionalidad mediante una API
REST desarrollada con FastAPI, que actúa como punto de entrada y salida del
flujo completo. El análisis se inicia a través de un endpoint que recibe el
archivo de logs a procesar; una vez iniciado, el avance del grafo se
transmite en tiempo real mediante \textit{Server-Sent Events}, lo que
permite a la interfaz reflejar el estado del análisis —qué nodo se
encuentra en ejecución, si se activó el ciclo de corrección, si el sistema
está consultando la memoria de largo plazo— mientras este ocurre, en lugar
de exponer únicamente un resultado final una vez concluido el
procesamiento. El resultado del análisis corresponde al reporte curado por
el Agente Auditor Final, que se persiste y queda disponible tanto en
formato estructurado como en PDF (\S\ref{sec:9.6-reportes}). Complementa
este flujo un punto de entrada alternativo por línea de comandos, utilizado
específicamente para las corridas de evaluación experimental descritas en
el Capítulo 10, donde la interacción mediante interfaz web resultaría poco
práctica dado el volumen de análisis ejecutados.

El Cuadro~\ref{tab:9.1-cap8-vs-cap9} sistematiza esta comparación,
consolidando en una sola tabla la correspondencia —o divergencia— entre
cada elemento de la propuesta del Capítulo 8 y su implementación efectiva,
como referencia para el resto del capítulo.

\begin{table}[H]
\centering
\small
\begin{tabular}{p{3.3cm}p{3cm}p{7cm}}
\toprule
\textbf{Elemento de la propuesta} & \textbf{Estado} & \textbf{Descripción} \\
\midrule
Arquitectura Plan-and-Execute (7 nodos, 3 \textit{routers}) & Implementado sin modificaciones & La topología del grafo construido coincide exactamente con la especificada en la Figura 8.3. \\
Modelo de lenguaje & Modificado respecto del diseño & El modelo especificado en la propuesta (\texttt{gemini-2.5-pro}) no estuvo disponible para la cuenta de desarrollo utilizada, por una restricción de acceso impuesta por el proveedor. Se adoptó \texttt{gemini-3.1-pro-preview} como modelo principal, con \texttt{gemini-3.6-flash} como respaldo automático. \\
Preprocesamiento de logs & Modificado y ampliado & Se implementó un patrón de adaptadores por formato (motor activo por defecto), junto con una segunda tubería de preprocesamiento heterogéneo de catorce fases, no contemplada en la propuesta original. \\
Memoria de largo plazo & Implementado y ampliado & Se incorporó una cuarta colección, dedicada a la base de conocimiento de MITRE ATT\&CK, además de las tres colecciones originalmente previstas. \\
Recuperación aumentada por generación & Implementado & Se definieron umbrales de similitud independientes por colección. \\
Integración de MITRE ATT\&CK (STIX 2.1) & Implementado & Se incorporó como una extensión de la herramienta de recuperación de conocimiento existente, sin necesidad de agregar un octavo nodo al grafo; este componente no se encontraba resuelto en el Capítulo 8. \\
Consulta de vulnerabilidades (NVD) & Implementado & La búsqueda se realiza a partir del primer servicio identificado en el resumen del lote de logs. \\
Reputación de direcciones IP (AbuseIPDB) & Implementado sin modificaciones & \\
Ciclo de corrección de comandos de mitigación & Implementado & Verificado mediante una suite de pruebas automatizadas. \\
Generación de reportes & Implementado y ampliado & Se incorporó la exportación a formato PDF, no contemplada en la propuesta. \\
Interfaz de usuario & Implementado y ampliado & La propuesta no contemplaba una interfaz; se desarrollaron dos: un panel de respaldo en HTML/JavaScript y una aplicación basada en Next.js orientada a la presentación del sistema. \\
Diseño experimental & Parcialmente implementado & Ejecutado sobre los conjuntos HDFS y BGL; el escenario de red basado en UWF-ZeekDataFall22 y la evaluación cualitativa mediante rúbrica de expertos permanecen pendientes (véase el Capítulo 10). \\
\bottomrule
\end{tabular}
\caption{Correspondencia entre la propuesta del Capítulo 8 y la implementación realizada.}
\label{tab:9.1-cap8-vs-cap9}
\end{table}

\section{Procesamiento de Datos}

\section{Implementación del Sistema Multiagente}

\subsection{Agentes LLM}

\subsection{Grafo de Estados y Routers}

\subsection{Estado Compartido del Sistema}

\section{Memoria y Recuperación de Información}

\subsection{Memoria de Corto Plazo}

\subsection{Memoria de Largo Plazo}

\subsection{Base de Datos Vectorial}

\subsection{Recuperación Aumentada por Generación}

\section{Integración de Herramientas}

\subsection{NVD y CVE}

\subsection{AbuseIPDB}

\subsection{MITRE ATT\&CK y STIX 2.1}

\section{Generación de Reportes}

\section{Interfaz y Despliegue}

\section{Pruebas de Implementación}